% Part: first-order-logic
% Chapter: models-theories
% Section: expressing-relations

\documentclass[../../../include/open-logic-section]{subfiles}

\begin{document}

% SEGMENT OLP-0171-S01

\olfileid{fol}{mat}{exr}

\olsection{\article{structure}
  \printtoken{S}{structure} இல் தொடர்புகளை வெளிப்படுத்துதல்}

\begin{explain}
!!{formula}s இன் முக்கியப் பயன்பாடுகளில் ஒன்று, !!a{structure}~$\Struct M$
இல் உள்ள பண்புகளையும் தொடர்புகளையும்,~$\Struct M$ இன் மொழி~$\Lang L$
இன் அடிப்படைக் குறியீடுகளைக் கொண்டு வெளிப்படுத்துதல். இதன் பொருள்
பின்வருவது. $\Struct M$ இன் !!{domain} ஒரு பொருட் தொகுப்பு.
!!{constant}s, !!{function}s மற்றும் !!{predicate}s ஆகியவை~$\Struct M$
இல் முறையே~$\Domain M$ இன் சில பொருட்கள்,~$\Domain M$ மீதான சார்புகள்,
~$\Domain M$ மீதான தொடர்புகள் ஆகியவற்றால் பொருள்கொள்ளப்படுகின்றன.
எடுத்துக்காட்டாக, $\Obj A^2_0$ ஆனது $\Lang L$ இல் இருந்தால்,
$\Struct M$ அதற்கு தொடர்பு~$R = \Assign{{\Obj A^2_0}}{M}$ ஐ ஒதுக்குகிறது.
அப்போது வாய்பாடு $\Atom{\Obj A^2_0}{\Obj v_1, \Obj v_2}$ பின்வரும்
பொருளில் அதே தொடர்பை \emph{வெளிப்படுத்துகிறது}: ஒரு மாறி ஒதுக்கீடு~$s$
ஆனது $\Obj v_1$ ஐ $a \in \Domain{M}$ க்கும் $\Obj v_2$ ஐ
$b \in \Domain M$ க்கும் கொண்டு சென்றால்,
\[
Rab \text{\quad iff\quad} \Sat{M}{\Atom{\Obj A^2_0}{\Obj v_1, \Obj v_2}}[s].
\]
இங்கு மாறி ஒதுக்கீடுகளைப் பயன்படுத்தியே ஆக வேண்டும் என்பதை கவனிக்க.
``$Rab$ iff $\Sat{M}{\Atom{\Obj A^2_0}{a, b}}$'' என்று மட்டும்
கூற முடியாது; ஏனெனில் $a$ மற்றும் $b$ நமது மொழியின் குறியீடுகள் அல்ல:
அவை~$\Domain{M}$ இன் !!{element}s.

நம்மிடம் அணு !!{formula}s மட்டும் இல்லை; அவற்றைத் தருக்க இணைப்பிகளாலும்
அளவையீட்டுக்குறிகளாலும் இணைக்கலாம். எனவே மேலும் சிக்கலான !!{formula}s,
~$\Struct M$ இல் நேரடியாக அமைக்கப்படாத பிற தொடர்புகளை வரையறுக்கலாம்.
இதை எவ்வாறு செய்வது என்பதிலும், குறிப்பாக !!a{structure} ஒன்றில் எந்தத்
தொடர்புகளை வரையறுக்க முடியும் என்பதிலும் நாம் ஆர்வம் கொண்டுள்ளோம்.
\end{explain}

\begin{defn}
$!A(\Obj v_1,\dots, \Obj v_n)$ ஆனது $\Obj v_1$,\dots, $\Obj v_n$
மட்டுமே கட்டற்றதாகத் தோன்றும் $\Lang L$ இன் !!a{formula} ஆகவும்,
$\Struct M$ ஆனது~$\Lang L$ க்கான !!a{structure} ஆகவும் இருக்கட்டும்.
$s(\Obj v_i) = a_i$ ($i = 1, \dots, n$) ஆகும் எந்த மாறி
ஒதுக்கீடு~$s$ க்கும் பின்வருவது பொருந்தினால்,
$!A(\Obj v_1,\dots, \Obj v_n)$ ஆனது
$R \subseteq \Domain M^n$ என்ற \emph{தொடர்பை வெளிப்படுத்துகிறது}:
\[
Ra_1\dots a_n \text{\quad iff\quad} \Sat{M}{\Atom{!A}{\Obj
    v_1,\dots, \Obj v_n}}[s]
\]
\end{defn}

\begin{ex}
எண்கணிதத்தின் நிலையான மாதிரி~$\Struct N$ இல், !!{formula}
$\Obj v_1 < \Obj v_2 \lor \eq[\Obj v_1][\Obj v_2]$ ஆனது~$\Nat$ மீதான
$\le$ தொடர்பை வெளிப்படுத்துகிறது. !!{formula}
$\eq[\Obj v_2][\Obj v_1']$ தொடர்வுறுப்புத் தொடர்பை வெளிப்படுத்துகிறது;
அதாவது, $m$ ஆனது~$n$ இன் தொடர்வுறுப்பு எனில் $Rnm$ பொருந்தும்
$R \subseteq \Nat^2$ என்ற தொடர்பை. வாய்பாடு
$\eq[\Obj v_1][\Obj v_2']$ முன்னுறுப்புத் தொடர்பை வெளிப்படுத்துகிறது.
!!{formula}s $\lexists[\Obj v_3][(\eq/[\Obj v_3][\Obj 0] \land
  \eq[\Obj v_2][(\Obj v_1 + \Obj v_3)])]$ மற்றும் $\lexists[\Obj
  v_3][\eq[(\Obj v_1 + {\Obj v_3}')][v_2]]$ இரண்டுமே $\Obj <$
தொடர்பை வெளிப்படுத்துகின்றன. இதனால் பயனிலைக் குறியீடு~$<$ உண்மையில்
எண்கணித மொழியில் தேவையற்றது; அதை வரையறுக்க முடியும்.
\end{ex}

\begin{explain}
இக்கருத்து குறிப்பிட்ட !!{structure}s இல் மட்டும் ஆர்வமூட்டுவதல்ல.
ஒரு மொழியைக் கொண்டு நோக்கமாகக் கொண்ட ஒரு மாதிரியையோ பல மாதிரிகளையோ
விவரிக்கும்போதெல்லாம், அதாவது கோட்பாடுகளைக் கருதும்போதெல்லாம், பொதுவாகவும்
ஆர்வமூட்டுகிறது. இக்கோட்பாடுகள் பெரும்பாலும் சில !!{predicate}s ஐ மட்டுமே
அடிப்படைக் குறியீடுகளாகக் கொண்டுள்ளன; ஆனால் அவை விவரிக்கும் பொருட்களத்தில்
வேறு பல தொடர்புகள் முக்கியப் பங்கு வகிக்கின்றன. மொழியின் அடிப்படை
!!{predicate}s ஐப் பொருள்கொள்ளும் தொடர்புகளைக் கொண்டு இப்பிற தொடர்புகளை
முறையாக வெளிப்படுத்த முடிந்தால், அவற்றை மொழியில் \emph{வரையறுக்க}
முடியும் என்கிறோம்.
\end{explain}

\begin{prob}
பின்வரும் தொடர்புகளை வரையறுக்கும் $\Lang L_A$ இன் !!{formula}s ஐக் காண்க:
\begin{enumerate}
\item $n$ ஆனது $i$ மற்றும் $j$ க்கு இடையில் உள்ளது;
\item $n$ ஆனது $m$ ஐ மீதமின்றி வகுக்கிறது (அதாவது, $m$ ஆனது $n$ இன்
  மடங்கு);
\item $n$ ஒரு பகா எண் (அதாவது, $1$ மற்றும் $n$ தவிர வேறு எந்த எண்ணும்
  ~ $n$ ஐ மீதமின்றி வகுக்காது).
\end{enumerate}
\end{prob}

\begin{prob}
வாய்பாடு $!A(\Obj v_1, \Obj v_2)$ ஆனது !!a{structure}~$\Struct M$
இல் தொடர்பு $R \subseteq \Domain M^2$ ஐ வெளிப்படுத்துகிறது என்க.
பின்வரும் தொடர்புகளை வெளிப்படுத்தும் வாய்பாடுகளைக் காண்க:
\begin{enumerate}
\item $R$ இன் நேர்மாறு $R^{-1}$;
\item சார்புப் பெருக்கல் $R \mid R$;
\end{enumerate}
$R$ இன் கடப்பு மூடல் $R^+$ ஐ வெளிப்படுத்த ஒரு வழியைக் காண முடியுமா?
\end{prob}

\begin{prob}
இரு-இடப் பயனிலைக் குறியீடு $<$ ஐ மட்டுமே கொண்ட மொழி $\Lang{L}$ ஆக
இருக்கட்டும் (வேறு !!{constant}s, !!{function}s அல்லது !!{predicate}s
இல்லை---நிச்சயமாக~$\eq$ மட்டும் உண்டு). பின்வருமாறு அமையும் கட்டமைப்பு
$\Struct{N}$ ஐக் கொள்க:
$\Domain{N} = \Nat$, மேலும் $\Assign{<}{N} = \Setabs{\tuple{n,m}}{n <
  m}$. பின்வருவனவற்றை நிரூபிக்க:
\begin{enumerate}
\item $\{ 0 \}$ ஆனது $\Struct{N}$ இல் வரையறுக்கத்தக்கது;
\item $\{ 1 \}$ ஆனது $\Struct{N}$ இல் வரையறுக்கத்தக்கது;
\item $\{ 2 \}$ ஆனது $\Struct{N}$ இல் வரையறுக்கத்தக்கது;
\item ஒவ்வொரு $n \in \Nat$ க்கும், தொகுப்பு $\{ n \}$ ஆனது
  $\Struct{N}$ இல் வரையறுக்கத்தக்கது;
\item $\Domain{N}$ இன் ஒவ்வொரு இறுதிக்குட்பட்ட உட்கணமும்
  $\Struct{N}$ இல் வரையறுக்கத்தக்கது;
\item $\Domain{N}$ இன் ஒவ்வொரு இணை-இறுதியான உட்கணமும்
  $\Struct{N}$ இல் வரையறுக்கத்தக்கது ($X \subseteq \Nat$ இணை-இறுதியானது
  iff $\Nat \setminus X$ இறுதிக்குட்பட்டது).
\end{enumerate}
\end{prob}


\end{document}
